% !TeX root = ../../Book3.tex
\section{Perron 方法}
\label{b3:sec:1602}

一般区域没有可直接写出的核。Perron 方法从所有低于边界数据的次调和函数出发；圆盘上的调和替换迫使它们的上包络满足平均等式。

\subsection{次函数}
\label{b3:subsec:160201}

称连续次调和函数 \(v:D\to\mathbb R\) 为边值 \(\varphi\) 的\term[perron-cixanshu]{Perron 次函数}[Perron subfunction]，若对每个 \(\xi\in\partial D\) 有 \(\limsup_{z\to\xi}v(z)\le\varphi(\xi)\)。本章只采用连续的次函数；这个族已足以完成连续边值的构造。

常数 \(m=\min\varphi\) 是次函数。每个次函数由次调和最大值原理满足 \(v\le M=\max\varphi\)，所以取上确界不会在内部发散。连续性允许在任何内部闭圆盘上作 Poisson 替换。

\subsection{Perron 族}
\label{b3:subsec:160202}

把全部次函数组成的族记为 \(\mathcal P_\varphi\)，上包络记为 \(H_\varphi(z)=\sup_{v\in\mathcal P_\varphi}v(z)\)。有限最大值仍在族中。

若 \(B\Subset D\)，把 \(v\) 在 \(B\) 内换成同边值的 Poisson 解，外部保持 \(v\)，记为 \(v^B\)。比较原理给 \(v^B\ge v\)。该函数在圆周连续，且仍次调和；下文通过先比较原次函数、再比较替换部分给出跨圆周的完整核查。替换只影响内部紧集，所以原边界限制不变。这个封闭性是上包络证明的核心。

\subsection{Perron 解}
\label{b3:subsec:160203}

\begin{theorem}[Perron]\label{b3:thm:16-perron}
有界区域 \(D\) 与连续边值 \(\varphi\) 的包络 \(H_\varphi\) 在 \(D\) 调和，并满足 \(m\le H_\varphi\le M\)。
\end{theorem}
\begin{proof}
固定内部圆盘 \(B\) 及 \(a\in B\)。选 \(v_n\) 使 \(v_n(a)\to H_\varphi(a)\)，逐次取有限最大值，得到递增序列，再在 \(B\) 调和替换。所得 \(h_n\) 递增、有界，在 \(B\) 局部一致收敛到调和函数 \(h\)，且 \(h(a)=H_\varphi(a)\)、\(h\le H_\varphi\)。
任取另一项 \(v\in\mathcal P_\varphi\)，将 \(\max(v,h_n)\)（在全域使用对应替换后的函数）再提升，得到调和函数 \(k_n\ge h_n,v\)。它们同样递增有界，极限 \(k\) 调和。中心满足 \(h(a)\le k(a)\le H_\varphi(a)=h(a)\)，所以非负调和函数 \(k-h\) 在内部为零，由强最小值原理恒零。于是 \(v\le h\) 于 \(B\)，取全族上确界得 \(H_\varphi\le h\)。因此包络在每个内部圆盘调和。
\end{proof}
证明没有先断言任意上包络连续；连续性由局部调和识别产生。这一步避免了把上确界与上半连续正则化混同。

\subsection{唯一性}
\label{b3:subsec:160204}

若一个经典解 \(u\) 已存在，则它本身是次函数，每个次函数又由比较原理不超过 \(u\)，故 \(H_\varphi=u\)。因此 Perron 构造不产生第二种不同的经典解。

没有边界正则性时，\(H_\varphi\) 仍是内部调和函数，却可能忽略某些指定的边界值。去心圆盘上，外圆取零、孤立边界点原点取一，得到 Perron 解恒零：有界调和函数的孤立奇点可去，延拓后外圆边值迫使零函数。该例将在下一节同时说明障碍函数为何必要。


\ProseParagraphBreak

调和替换在圆周上的粘合可以作如下严格核查。设 \(w=v^B\)，并在测试圆盘 \(Q\Subset D\) 取调和函数 \(h\)，满足 \(w\le h\) 于 \(\partial Q\)。由于 \(v\le w\)，先由 \(v\) 的次调和比较得到 \(v\le h\) 于 \(Q\)。在 \(Q\cap B\) 中，\(w-h\) 调和；其边界分成 \(\partial Q\cap\overline B\) 与 \(Q\cap\partial B\)，前者由假设非正，后者因 \(w=v\le h\) 非正。最大值原理便给 \(w\le h\) 于交集。其余部分 \(w=v\le h\)，所以整个 \(Q\) 比较成立。这里先比较原次函数，再比较替换部分，避免在交界处循环使用尚未证明的次调和性。

\ProseParagraphBreak

Perron 证明中选择序列只需要在一个点逼近包络，而不要求全域同时选出最优函数。具体可令 \(v_n(a)>H_\varphi(a)-1/n\)，将 \(m,v_1,\ldots,v_n\) 的最大值记为 \(s_n\)，再替换得 \(w_n=s_n^B\)。序列 \(w_n\) 在全域递增，且属于同一 Perron 族；其圆盘限制 \(h_n\) 调和。由于 \(h_n(a)\ge s_n(a)\) 且 \(h_n\le H_\varphi\)，中心极限正是包络。对任意竞争者 \(v\)，使用 \(\max(v,w_n)\) 而非只定义在 \(B\) 中的 \(\max(v,h_n)\)，就得到合法的全域次函数。这些域的区别保证第二轮替换仍保留边界约束。

\ProseParagraphBreak

取上确界并不通常与取极限交换。Perron 证明采用的事实是：对每个固定竞争者 \(v\)，新调和极限 \(k\) 在圆盘上支配它，同时 \(k(a)=h(a)\)；非负调和函数的强最小值原理强迫 \(k=h\)。于是一个最初只在中心被优化的序列，最终在整个圆盘支配全部竞争者。这个步骤使用线性的调和方程和正性，才产生包络的调和性；只看有限最大值的次调和封闭性还不足以完成证明。

\ProseParagraphBreak

去心圆盘例子的包络恒零也能在不预先调用可去奇点定理的情况下证明。对任意次函数 \(v\)，外圆边界上限零，内边界点上限至多一，故 \(v\le1\)。给定 \(\varepsilon>0\)，在足够小的内圆 \(|z|=r\) 上 \(v\le1+\varepsilon\)，环域上的调和函数
\[
q_r(z)=(1+\varepsilon)\frac{\log|z|}{\log r}
\]
外圆值零、内圆值 \(1+\varepsilon\)。比较给 \(v(z)\le q_r(z)\)；固定 \(z\ne0\) 让 \(r\downarrow0\)，得到 \(v(z)\le0\)。常数零又是合法次函数，所以包络恰为零。孤立边值一在内部没有影响，这一计算完全在当前已证明的环域比较范围内进行。

\ProseParagraphBreak

边界数据的单调性不需要域正则：若 \(\varphi\le\psi\)，则 \(\mathcal P_\varphi\subseteq\mathcal P_\psi\)，故 \(H_\varphi\le H_\psi\)。加常数与乘正数也分别对应次函数族的平移与缩放，所以 \(H_{\varphi+c}=H_\varphi+c\)、\(H_{a\varphi}=aH_\varphi\) 对 \(a>0\) 成立。两者推出
\[
\|H_\varphi-H_\psi\|_{L^\infty(D)}\le\|\varphi-\psi\|_{L^\infty(\partial D)}.
\]
一般加法的线性不能只由族包含轻率推出；本章在正则域用经典解的唯一性证明完整线性，已经足够用于测度表示。

\ProseParagraphBreak

圆环提供检验一般存在理论的另一组显式解。设 \(r_0<|z|<R\)，令 \(x=\log(|z|/r_0)\)、\(L=\log(R/r_0)\)，则 \(0<x<L\)。极坐标下 Laplace 方程乘以 \(|z|^2\) 后成为 \(u_{xx}+u_{\theta\theta}=0\)。对内、外边界各一个 Fourier 模式，解分别为
\[
 \frac{\sinh(n(L-x))}{\sinh(nL)}e^{\mathrm{i}n\theta},\qquad
 \frac{\sinh(nx)}{\sinh(nL)}e^{\mathrm{i}n\theta}\quad(n\ge1).
\]
它们在相应边界取一倍模式，在另一边界取零；对负频率取共轭。零频率则为 \(1-x/L\) 与 \(x/L\)。每一式都可直接求导检验，唯一性保证它与 Perron 解相同。
\ProseParagraphBreak
若内、外边界为三角多项式，有限叠加便得到解。对一般连续周期数据，可由非负的 Fej\'er 核逼近：
\[
 F_N(t)=\frac1N\left|\sum_{j=0}^{N-1}e^{\mathrm{i}jt}\right|^2,
 \qquad \frac1{2\pi}\int_{-\pi}^{\pi}F_N(t)\,\mathrm{d}t=1.
\]
在 \(|t|\ge\delta\) 上，有限几何和给 \(F_N(t)\le C_\delta/N\)；在小弧上用数据的一致连续性。因此卷积 \(F_N*\varphi\) 一致趋于 \(\varphi\)，每次卷积都是三角多项式。对这些数据的圆环解，最大值原理使内部一致差不超过两条边界数据的最大差，故在闭圆环上一致收敛。极限在内部调和，因为任意小圆盘上的平均值性质可通过一致极限；边值也正确。这给出圆环上存在性的第二个完整构造，并显示对数项恰来自角向零频率。
